\documentclass[11pt]{article}
% General SciPost-style report. Use the official journal class when required for submission.
\usepackage[a4paper,margin=22mm]{geometry}
\usepackage{amsmath,amssymb,amsthm,mathtools}
\usepackage{graphicx,booktabs,enumitem,hyperref,xcolor}
\newtheorem{theorem}{Theorem}
\newtheorem{lemma}[theorem]{Lemma}
\newtheorem{proposition}[theorem]{Proposition}
\newtheorem{definition}[theorem]{Definition}
\title{NIMA_REPORT_TITLE}
\author{Author names}
\date{\today}
\begin{document}
\maketitle
\begin{abstract}
State the question, actual result, assumptions and remaining uncertainty.
\end{abstract}

NIMA_REPORT_BODY

% Use vector PDF plots or hybrid illustrations at their designed physical size.
% Place detailed figure explanations from the companion Markdown in captions or surrounding prose.
% Example: \includegraphics[width=\linewidth]{../../figures/example/figure.pdf}
\appendix
\section{Vocabulary and conventions}
Define the symbols, domains and project-specific terms used in this document.
Identify the sealed vocabulary and its revision when available; distinguish a provisional glossary.

\section{Evidence, provenance and reproducibility}
Record exact sources, calculation or proof sessions, code, data and execution references.
Put detailed software identifiers here when they would interrupt the scientific explanation.

\section{Open obligations and limitations}
Report every open obligation relevant to the reported results and retain additional audit-only obligations separately.
Distinguish completed operations from established physical or mathematical conclusions.
\end{document}
